\documentclass[aspectratio=169, 11pt, xcolor={table}]{beamer}
% Shared style for the algorithmics course slides (same template as the
% ENSB generative-AI course), with C code listings added.
% Dark blue removed: every slide (cover included) uses the light background
% of the template's content slides, with teal + orange as accents.

\usepackage{amsmath}
\usepackage{amssymb}
\usepackage{booktabs}
\usepackage{colortbl}
\usepackage{tabularx}
\usepackage{tikz}
\usepackage[skins]{tcolorbox}
\usepackage{fontawesome5}
\usepackage{listings}
\usepackage{polyglossia}

\usetikzlibrary{positioning, calc, arrows.meta, shapes.geometric}

%------------------------------------------------
% Languages and fonts (compile from slides/ with XeLaTeX or tectonic)
%------------------------------------------------
\setdefaultlanguage[numerals=maghrib]{arabic}
\setotherlanguage{english}
\usefonttheme{serif}

\newfontfamily\arabicfont[
    Script=Arabic, Scale=0.98, Path=../../fonts/, Extension=.ttf,
    UprightFont=*-Regular, BoldFont=*-Bold,
    ItalicFont=*-Italic, SlantedFont=*-Italic,
    BoldItalicFont=*-BoldItalic, BoldSlantedFont=*-BoldItalic]{Amiri}
\newfontfamily\arabicfontsf[
    Script=Arabic, Scale=0.98, Path=../../fonts/, Extension=.ttf,
    UprightFont=*-Regular, BoldFont=*-Bold,
    ItalicFont=*-Italic, SlantedFont=*-Italic,
    BoldItalicFont=*-BoldItalic, BoldSlantedFont=*-BoldItalic]{Amiri}
% Arabic fallback for \ttfamily (needed by the C listings)
\newfontfamily\arabicfonttt[Script=Arabic, Scale=0.98, Path=../../fonts/,
    Extension=.ttf, UprightFont=*-Regular, BoldFont=*-Bold,
    ItalicFont=*-Italic, SlantedFont=*-Italic,
    BoldItalicFont=*-BoldItalic, BoldSlantedFont=*-BoldItalic]{Amiri}
% Latin monospace font for C code (selected directly, so polyglossia
% does not swap it for the Arabic fallback above)
\newfontfamily\codefont[Scale=0.95, BoldFont=lmmonolt10-bold.otf,
    ItalicFont=lmmono10-italic.otf]{lmmono10-regular.otf}
% Heading font
\newfontfamily\khatt[Script=Arabic, Scale=1.05, Path=../../fonts/,
    BoldFont=ae_AlBattar.ttf]{ae_AlBattar.ttf}

%------------------------------------------------
% Colors (template palette without the navy tones)
%------------------------------------------------
\definecolor{brandTeal}{HTML}{5CC8E6}
\definecolor{brandOrange}{HTML}{FFA27F}
\definecolor{brandLight}{HTML}{F8F9FA}
\definecolor{brandDark}{HTML}{263238}
\definecolor{brandMuted}{HTML}{6C757D}
\colorlet{tealDeep}{brandTeal!60!black}
\definecolor{codeGrey}{HTML}{F1F3F5}

\setbeamercolor{background canvas}{bg=brandLight}
\setbeamercolor{normal text}{fg=brandDark}
\setbeamercolor{structure}{fg=brandTeal}
\setbeamertemplate{navigation symbols}{}

% beamer.cls sets \raggedright globally; under RTL that centers short lines
% or pushes them left, so go back to justified text (flush right).
\rightskip=0pt\relax
% Lists call \raggedright too; make it a no-op for RTL slides.
\AtBeginDocument{\renewcommand{\raggedright}{\rightskip=0pt\relax}}

%------------------------------------------------
% Frame title, footer, lists
%------------------------------------------------
\setbeamertemplate{frametitle}{%
    \vspace{0.35cm}
    \noindent{\khatt\bfseries\Large\color{brandDark}\insertframetitle}
    \par
    \vskip-2pt
    \noindent{\color{brandTeal}\rule{\textwidth}{1.5pt}}
    \vspace{0.05cm}
}

% Footer: page number only
\setbeamertemplate{footline}{%
    \hbox to \paperwidth{\hfill
        {\color{brandMuted}\scriptsize\LR{\insertframenumber\,/\,\inserttotalframenumber}}%
        \hspace{0.6cm}}%
    \vspace{0.25cm}}

\setbeamertemplate{itemize item}{\color{brandTeal}\small\faChevronLeft}
\setbeamertemplate{itemize subitem}{\color{brandMuted}\small\faAngleLeft}
\setbeamertemplate{enumerate item}{\color{tealDeep}\bfseries\insertenumlabel)}

% Tables: every cell framed by a visible border
\newcolumntype{R}{>{\raggedleft\arraybackslash}X}
\newcolumntype{C}{>{\centering\arraybackslash}X}
\newcommand{\ras}[1]{\cellcolor{brandTeal!20}\textbf{#1}}
% Colored text inside a table cell (\leavevmode keeps it on the row baseline)
\newcommand{\cell}[2]{\leavevmode\textcolor{#1}{#2}}
\renewcommand{\arraystretch}{1.35}
\setlength{\arrayrulewidth}{0.6pt}
\arrayrulecolor{brandMuted!70}

%------------------------------------------------
% Boxes (the two boxes of the template's 2nd slide)
%------------------------------------------------
\tcbset{
    brandbox/.style={
        enhanced, arc=4pt, outer arc=4pt, boxrule=1.5pt,
        coltitle=white, coltext=brandDark, fonttitle=\khatt\bfseries,
        toptitle=0.8mm, bottomtitle=0.8mm,
        top=1.2mm, bottom=1.2mm, left=3mm, right=3mm,
        before skip=4pt, after skip=4pt}
}
% Teal box: definitions, objectives, summaries
\newtcolorbox{tbox}[2][]{brandbox, colback=brandTeal!10,
    colframe=brandTeal, title={#2}, #1}
% Orange box: warnings
\newtcolorbox{obox}[2][]{brandbox, colback=brandOrange!10,
    colframe=brandOrange, title={#2}, #1}
% Card (template's 3rd slide, in teal)
\newtcolorbox{card}[2][]{enhanced, colback=brandTeal!6, colframe=brandTeal,
    arc=6pt, boxrule=1.2pt, halign=center, halign title=center,
    coltitle=white, coltext=brandDark, fonttitle=\khatt\bfseries,
    top=1.5mm, bottom=1.5mm, left=2mm, right=2mm, title={#2}, #1}
% Prompt text written to the tool
\newtcolorbox{amr}[1][brandTeal]{enhanced, arc=1mm, boxrule=0pt,
    colback=white, overlay={\fill[#1] (frame.north east) rectangle ([xshift=-2.5pt]frame.south east);},
    left=6pt, right=8pt, top=2pt, bottom=2pt,
    before skip=3pt, after skip=3pt}

%------------------------------------------------
% TikZ (Arabic right-to-left inside nodes)
%------------------------------------------------
\newcommand{\nodecolor}{}
\tikzset{
    every node/.append style={
        execute at begin node={\nodecolor\everypar{\beginR}\beginR},
        execute at end node=\endR},
    lawn/.code={\def\nodecolor{\color{#1}}},
    khatwa/.style={draw=#1, fill=#1!12, very thick, rounded corners=5pt,
        minimum height=0.95cm, text width=1.9cm, align=center,
        font=\khatt\bfseries\small},
    khatwa/.default=brandTeal,
    sahm/.style={-{Stealth[length=2.5mm]}, thick, color=brandMuted},
    rang/.style={draw=#1, fill=#1!12},
    raqm/.style={circle, fill=#1, lawn=white, inner sep=1.5pt,
        font=\small\bfseries},
}

% Horizontal chain of steps, right to left:
% \flow[width]{{a}, {b\\c}, {d}}
\newcommand{\flow}[2][1.9cm]{%
    \begin{tikzpicture}[node distance=5mm]
        \foreach \t [count=\i] in {#2} {
            \ifnum\i=1
                \node[khatwa, text width=#1] (f\i) {\t};
            \else
                \pgfmathtruncatemacro\p{\i-1}
                \node[khatwa, text width=#1, left=of f\p] (f\i) {\t};
                \draw[sahm] (f\p) -- (f\i);
            \fi
        }
    \end{tikzpicture}}

%------------------------------------------------
% C code (always left-to-right, inside a fragile frame)
%------------------------------------------------
\lstdefinestyle{cstyle}{
    language=C, basicstyle=\codefont\footnotesize, columns=fullflexible,
    keepspaces=true, showstringspaces=false, upquote=true,
    keywordstyle=\color{tealDeep}\bfseries,
    commentstyle=\color{red!75!black},
    stringstyle=\color{brandOrange!85!black},
    xleftmargin=16pt,
    numbers=left, numberstyle=\tiny\color{brandMuted}, numbersep=7pt,
    aboveskip=0pt, belowskip=0pt,
    morekeywords={bool, true, false}}
% Inline listings must sit in the english environment, not LTR, otherwise
% numbers such as 3.14 are reversed.
% Numbers in Arabic text: \textenglish{1.75} (not \LR).
% Compact code size for long programs (fixed line height, no gaps)
\newcommand{\codesmall}{\fontsize{8}{8}\selectfont}
\newcommand{\codetiny}{\fontsize{7.4}{7.6}\selectfont}
% Inline C: \code{printf}
\newcommand{\code}[1]{\mbox{\textenglish{\codefont\begingroup\color{tealDeep}#1\endgroup}}}
% Program output (console)
\lstdefinestyle{outstyle}{
    language={}, basicstyle=\codefont\footnotesize\color{brandDark},
    columns=fullflexible, keepspaces=true, showstringspaces=false,
    aboveskip=0pt, belowskip=0pt}

% Listings read from files (no fragile frame needed):
%   \ccode{d1/circle.c}   \out{d1/circle.out}
% One box per listing for the grey background (listings' own per-line
% background leaves thin white gaps between lines)
\newtcolorbox{codebox}{enhanced, colback=codeGrey, colframe=codeGrey,
    boxrule=0pt, arc=0pt, left=4pt, right=4pt, top=1pt, bottom=1pt,
    before skip=0pt, after skip=2pt,
    overlay={\fill[brandMuted!45] (frame.south west) rectangle ([xshift=2.5pt]frame.north west);}}
\newtcolorbox{outbox}{enhanced, colback=brandMuted!20, colframe=brandMuted!20,
    boxrule=0pt, arc=0pt, left=4pt, right=4pt, top=1.5pt, bottom=1.5pt,
    before skip=2pt, after skip=2pt}
\newcommand{\ccode}[2][]{\begin{codebox}\begin{english}\lstinputlisting[style=cstyle,#1]{../code/#2}\end{english}\end{codebox}}
\newcommand{\out}[2][]{\begin{outbox}\begin{english}\lstinputlisting[style=outstyle,#1]{../code/#2}\end{english}\end{outbox}}
% Audience line of the cover
\newcommand{\audience}{طلبة السنة الأولى}

%------------------------------------------------
% Cover (template's 1st slide on the light background)
%------------------------------------------------
\newcommand{\cover}[3]{%
\begin{frame}[plain]
    \begin{tikzpicture}[remember picture, overlay]
        \draw[brandTeal, line width=1.5pt, opacity=0.3]
            (current page.south east) ++(-1,1) circle (4cm);
        \draw[brandTeal, line width=0.8pt, opacity=0.25]
            (current page.north west) ++(2,2) circle (6cm);
        \draw[brandOrange, line width=0.5pt, opacity=0.25]
            (current page.north west) ++(2,2) circle (7cm);
        \node[anchor=north east] at ($(current page.north east)+(-1.0,-0.8)$) {
            \begin{tikzpicture}[scale=0.35]
                \fill[brandTeal] (0,0) -- (0.8,1.2) -- (1.6,0) -- cycle;
                \fill[brandOrange] (0.8,1.2) -- (1.6,2.4) -- (2.4,1.2) -- cycle;
            \end{tikzpicture}
        };
    \end{tikzpicture}

    \vspace{1.6cm}
    \begin{minipage}{0.85\textwidth}
        {\color{tealDeep}\khatt\bfseries #1} \\[0.3cm]
        {\raggedleft\color{brandDark}\Huge\khatt\bfseries #2\par} \vspace{0.3cm}
        {\color{brandMuted}\large #3}
    \end{minipage}

    \vspace{1.2cm}
    \begin{minipage}{0.85\textwidth}
        \color{brandMuted}\small
        \begin{tabular}{@{}ll}
            \textbf{الأستاذ:} & د. بن عبد الرزاق يوسف \\
            \textbf{الفئة:}   & \audience \\
        \end{tabular}
    \end{minipage}
\end{frame}}

\title{العودية}

\begin{document}

\cover{درس الفصل الثامن}
{العودية}
{الدوال التي تستدعي نفسها: حالة التوقف والاستدعاء العودي}

\begin{frame}{أهداف الدرس}
    \begin{tbox}{\faBullseye\quad في نهاية هذا الدرس، تستطيع أن:}
        \begin{itemize}
            \item تعرّف الدالة العودية وتحدد حالة التوقف.
            \item تتتبع الاستدعاءات العودية في المكدس.
            \item تقارن بين الحل العودي والحل التكراري.
        \end{itemize}
    \end{tbox}
\end{frame}

\begin{frame}{الدالة العودية}
    \begin{tbox}{\faLightbulb\quad تعريف}
        \textbf{الدالة العودية}: دالة تستدعي \textbf{نفسها} لحل نسخة أصغر من المشكلة نفسها.
    \end{tbox}
\end{frame}

\begin{frame}{جزءا الدالة العودية}
    \vspace{0.6cm}
    \begin{columns}[T]
        \column{0.45\textwidth}
        \begin{card}{حالة التوقف}
            حالة بسيطة نعرف حلها مباشرة، دون استدعاء.
        \end{card}
        \column{0.45\textwidth}
        \begin{card}{الاستدعاء العودي}
            نستدعي الدالة على مشكلة \textbf{أصغر}، تقترب من حالة التوقف.
        \end{card}
    \end{columns}
    \vspace{0.4cm}
    مثال: $n! = n \times (n-1)!$ \quad و\quad $0! = 1$
\end{frame}

\begin{frame}[t]{العاملي بطريقة عودية}
    \begin{columns}[T]
        \column{0.63\textwidth}
        \ccode[basicstyle=\codefont\codesmall]{d8/l01.c}
        \column{0.34\textwidth}
        {\small كل استدعاء ينقص \code{n} بواحد حتى نصل إلى \code{0}.\par}
        \vspace{0.3cm}
        \out[basicstyle=\codefont\scriptsize\color{brandDark}]{d8/l01.out}
    \end{columns}
\end{frame}

\begin{frame}{تتبع \textenglish{fact(4)}}
    \begin{center}
    \begin{tikzpicture}[
        call/.style={draw=brandTeal, fill=brandTeal!12, thick, rounded corners=3pt,
            minimum width=4.2cm, minimum height=0.55cm, font=\small},
        stop/.style={call, draw=brandOrange, fill=brandOrange!15},
        val/.style={draw=brandOrange, fill=white, thick, circle, inner sep=1pt,
            minimum size=0.6cm, font=\small\bfseries}]
        \foreach \t/\v [count=\i from 0] in {{fact(4) = 4 * fact(3)}/24,
                {fact(3) = 3 * fact(2)}/6, {fact(2) = 2 * fact(1)}/2,
                {fact(1) = 1 * fact(0)}/1} {
            \node[call] (c\i) at (0, -0.85*\i) {\code{\t}};
            \node[val]  (v\i) at (3.3, -0.85*\i) {\textenglish{\v}};
        }
        \node[stop] (c4) at (0, -3.4) {\code{fact(0) = 1}};
        \node[val]  (v4) at (3.3, -3.4) {\textenglish{1}};
        \foreach \i [evaluate=\i as \j using int(\i+1)] in {0,...,3} {
            \draw[sahm] ([xshift=-1.6cm]c\i.south) -- ([xshift=-1.6cm]c\j.north);
            \draw[sahm, brandOrange] (v\j) -- (v\i);
        }
        \node[font=\scriptsize\bfseries, lawn=brandMuted] at (-3.3, -1.7) {النزول};
        \node[font=\scriptsize\bfseries, lawn=brandOrange!80!black] at (4.3, -1.7) {الصعود};
    \end{tikzpicture}
    \end{center}
    ننزل حتى حالة التوقف، ثم نصعد بالنتائج: \code{fact(4) = 24}.
\end{frame}

\begin{frame}[t]{مجموع الأعداد من 1 إلى n}
    \begin{columns}[T]
        \column{0.56\textwidth}
        \ccode[basicstyle=\codefont\codesmall]{d8/l02.c}
        \column{0.41\textwidth}
        {\small \code{sum(n) = n + sum(n-1)} و\code{sum(1) = 1}.\par}
        \vspace{0.3cm}
        \out[basicstyle=\codefont\scriptsize\color{brandDark}]{d8/l02.out}
    \end{columns}
\end{frame}

\begin{frame}[t]{القوة}
    \begin{columns}[T]
        \column{0.56\textwidth}
        \ccode[basicstyle=\codefont\codesmall]{d8/l03.c}
        \column{0.41\textwidth}
        {\small $x^n = x \times x^{n-1}$ و$x^0 = 1$.\par}
        \vspace{0.3cm}
        \out[basicstyle=\codefont\scriptsize\color{brandDark}]{d8/l03.out}
    \end{columns}
\end{frame}

\begin{frame}[t]{متتالية فيبوناتشي}
    \begin{columns}[T]
        \column{0.56\textwidth}
        \ccode[basicstyle=\codefont\codesmall]{d8/l04.c}
        \column{0.41\textwidth}
        {\small استدعاءان عوديان في كل مرة.\par}
        \vspace{0.3cm}
        \out[basicstyle=\codefont\scriptsize\color{brandDark}]{d8/l04.out}
    \end{columns}
\end{frame}

\begin{frame}{شجرة استدعاءات \textenglish{fib(4)}}
    \begin{center}
    \begin{tikzpicture}[
        level distance=1.05cm,
        level 1/.style={sibling distance=4.6cm},
        level 2/.style={sibling distance=2.3cm},
        level 3/.style={sibling distance=1.2cm},
        every node/.append style={draw=brandTeal, fill=brandTeal!12, thick,
            rounded corners=3pt, font=\scriptsize, inner sep=3pt},
        tikrar/.style={draw=brandOrange, fill=brandOrange!20},
        edge from parent/.style={draw, brandMuted, thick}]
        \node {\code{fib(4)}}
            child {node {\code{fib(3)}}
                child {node[tikrar] {\code{fib(2)}}
                    child {node {\code{fib(1)}}}
                    child {node {\code{fib(0)}}}}
                child {node {\code{fib(1)}}}}
            child {node[tikrar] {\code{fib(2)}}
                child {node {\code{fib(1)}}}
                child {node {\code{fib(0)}}}};
    \end{tikzpicture}
    \end{center}
    \code{fib(2)} يُحسب مرتين: الحل العودي هنا بطيء جدًا عندما يكبر \code{n}.
\end{frame}

\begin{frame}[t]{القاسم المشترك الأكبر}
    \begin{columns}[T]
        \column{0.61\textwidth}
        \ccode[basicstyle=\codefont\codesmall]{d8/l05.c}
        \column{0.36\textwidth}
        {\small خوارزمية إقليدس: \code{gcd(a, b) = gcd(b, a \% b)}.\par}
        \vspace{0.3cm}
        \out[basicstyle=\codefont\scriptsize\color{brandDark}]{d8/l05.out}
    \end{columns}
\end{frame}

\begin{frame}[t]{عرض أرقام عدد بالعكس}
    \begin{columns}[T]
        \column{0.56\textwidth}
        \ccode[basicstyle=\codefont\codesmall]{d8/l06.c}
        \column{0.41\textwidth}
        {\small دالة \code{void} عودية: \code{return} وحدها تكفي للتوقف.\par}
        \vspace{0.3cm}
        \out[basicstyle=\codefont\scriptsize\color{brandDark}]{d8/l06.out}
    \end{columns}
\end{frame}

\begin{frame}{دون حالة توقف}
    \begin{center}
    \begin{minipage}{0.7\textwidth}
        \ccode{d8/l07.c}
    \end{minipage}
    \end{center}
    \begin{obox}{\faExclamationTriangle\quad انتبه}
        تتراكم الاستدعاءات في المكدس حتى يتوقف البرنامج بخطأ \textenglish{stack overflow}.
    \end{obox}
\end{frame}

\begin{frame}{العودية أم التكرار؟}
    \begin{center}
    \begin{tabular}{|c|c|c|}
        \hline
        \ras{} & \ras{العودية} & \ras{التكرار (الحلقات)} \\ \hline
        الكتابة & قصيرة وقريبة من التعريف الرياضي & أطول أحيانًا \\ \hline
        الذاكرة & كل استدعاء يأخذ مكانًا في المكدس & ذاكرة ثابتة \\ \hline
        الاستعمال & الأشجار، التقسيم، البحث الثنائي & الحسابات البسيطة \\ \hline
    \end{tabular}
    \end{center}
\end{frame}

\begin{frame}{أمثلة تطبيقية}
    \begin{tbox}{\faLaptopCode\quad عشرة أمثلة من مجالات مختلفة}
        نطبق ما تعلمناه في مجالات من الحياة اليومية:
        الرياضة، الشبكات الاجتماعية، البيئة، ألعاب الفيديو، الشبكات، علم الأحياء، الفيزياء، محركات البحث، معالجة النصوص، التجارة الإلكترونية.
    \end{tbox}
\end{frame}

\begin{frame}[t]{المثال الأول: الرياضة}
    \begin{columns}[T]
        \column{0.61\textwidth}
        \ccode[basicstyle=\codefont\codesmall]{d8/ex01.c}
        \column{0.36\textwidth}
        {\small \textbf{عدد مباريات الإقصاء}: في كل دور تُلعب نصف عدد الفرق مباريات.\par}
        \vspace{0.3cm}
        \out[basicstyle=\codefont\scriptsize\color{brandDark}]{d8/ex01.out}
    \end{columns}
\end{frame}

\begin{frame}[t]{المثال الثاني: الشبكات الاجتماعية}
    \begin{columns}[T]
        \column{0.65\textwidth}
        \ccode[basicstyle=\codefont\codesmall]{d8/ex02.c}
        \column{0.32\textwidth}
        {\small \textbf{انتشار منشور}: كل شخص يشارك المنشور مع 3 أشخاص في المستوى التالي.\par}
        \vspace{0.3cm}
        \out[basicstyle=\codefont\scriptsize\color{brandDark}]{d8/ex02.out}
    \end{columns}
\end{frame}

\begin{frame}[t]{المثال الثالث: البيئة}
    \begin{columns}[T]
        \column{0.64\textwidth}
        \ccode[basicstyle=\codefont\codesmall]{d8/ex03.c}
        \column{0.33\textwidth}
        {\small \textbf{نمو غابة}: عدد الأشجار يزيد 20\% كل سنة، مع غرس 50 شجرة جديدة.\par}
        \vspace{0.3cm}
        \out[basicstyle=\codefont\scriptsize\color{brandDark}]{d8/ex03.out}
    \end{columns}
\end{frame}

\begin{frame}[t]{المثال الرابع: ألعاب الفيديو}
    \begin{columns}[T]
        \column{0.60\textwidth}
        \ccode[basicstyle=\codefont\codesmall]{d8/ex04.c}
        \column{0.37\textwidth}
        {\small \textbf{أبراج هانوي}: ننقل 3 أقراص من العمود \code{A} إلى العمود \code{C}.\par}
        \vspace{0.3cm}
        \out[basicstyle=\codefont\scriptsize\color{brandDark}]{d8/ex04.out}
    \end{columns}
\end{frame}

\begin{frame}[t]{المثال الخامس: الشبكات}
    \begin{columns}[T]
        \column{0.56\textwidth}
        \ccode[basicstyle=\codefont\codesmall]{d8/ex05.c}
        \column{0.41\textwidth}
        {\small \textbf{التحويل إلى الثنائي}: الحواسيب والشبكات تستعمل النظام الثنائي.\par}
        \vspace{0.3cm}
        \out[basicstyle=\codefont\scriptsize\color{brandDark}]{d8/ex05.out}
    \end{columns}
\end{frame}

\begin{frame}[t]{المثال السادس: علم الأحياء}
    \begin{columns}[T]
        \column{0.58\textwidth}
        \ccode[basicstyle=\codefont\codesmall]{d8/ex06.c}
        \column{0.39\textwidth}
        {\small \textbf{تكاثر الأرانب}: مسألة فيبوناتشي الأصلية: عدد أزواج الأرانب كل شهر.\par}
        \vspace{0.3cm}
        \out[basicstyle=\codefont\scriptsize\color{brandDark}]{d8/ex06.out}
    \end{columns}
\end{frame}

\begin{frame}[t]{المثال السابع: الفيزياء}
    \begin{columns}[T]
        \column{0.66\textwidth}
        \ccode[basicstyle=\codefont\codesmall]{d8/ex07.c}
        \column{0.31\textwidth}
        {\small \textbf{ارتداد كرة}: بعد كل ارتداد تفقد الكرة 30\% من ارتفاعها.\par}
        \vspace{0.3cm}
        \out[basicstyle=\codefont\scriptsize\color{brandDark}]{d8/ex07.out}
    \end{columns}
\end{frame}

\begin{frame}[t]{المثال الثامن: محركات البحث}
    \begin{columns}[T]
        \column{0.65\textwidth}
        \ccode[basicstyle=\codefont\codesmall]{d8/ex08.c}
        \column{0.32\textwidth}
        {\small \textbf{البحث الثنائي}: نبحث في جدول مرتب بتقسيمه إلى نصفين في كل مرة.\par}
        \vspace{0.3cm}
        \out[basicstyle=\codefont\scriptsize\color{brandDark}]{d8/ex08.out}
    \end{columns}
\end{frame}

\begin{frame}[t]{المثال التاسع: معالجة النصوص}
    \begin{columns}[T]
        \column{0.63\textwidth}
        \ccode[basicstyle=\codefont\codesmall]{d8/ex09.c}
        \column{0.34\textwidth}
        {\small \textbf{كلمة متناظرة}: نقارن الحرف الأول بالأخير، ثم نعيد على ما بينهما.\par}
        \vspace{0.3cm}
        \out[basicstyle=\codefont\scriptsize\color{brandDark}]{d8/ex09.out}
    \end{columns}
\end{frame}

\begin{frame}[t]{المثال العاشر: التجارة الإلكترونية}
    \begin{columns}[T]
        \column{0.60\textwidth}
        \ccode[basicstyle=\codefont\codesmall]{d8/ex10.c}
        \column{0.37\textwidth}
        {\small \textbf{مجموع السلة}: مجموع السلة = آخر منتج + مجموع باقي المنتجات.\par}
        \vspace{0.3cm}
        \out[basicstyle=\codefont\scriptsize\color{brandDark}]{d8/ex10.out}
    \end{columns}
\end{frame}

\begin{frame}{أخطاء شائعة}
    \begin{obox}{\faExclamationTriangle\quad تجنّب هذه الأخطاء}
        \begin{itemize}
            \item نسيان حالة التوقف، أو حالة توقف لا نصل إليها أبدًا.
            \item استدعاء عودي لا يصغّر المشكلة: \code{f(n)} تستدعي \code{f(n)}.
            \item استعمال العودية لحسابات تتكرر كثيرًا (فيبوناتشي) مع قيم كبيرة.
        \end{itemize}
    \end{obox}
\end{frame}

\begin{frame}{خلاصة الدرس}
    \begin{tbox}{\faCheckCircle\quad خلاصة الدرس}
        \begin{itemize}
            \item الدالة العودية تستدعي نفسها على مشكلة أصغر.
            \item كل دالة عودية تحتاج إلى حالة توقف.
            \item العودية مناسبة للأشجار والتقسيم، والتكرار أوفر في الذاكرة.
        \end{itemize}
    \end{tbox}
\end{frame}

\end{document}
